\section{Scope, Method, and Qualifications}

\subsection{Reader, question, and claim}

This is a discursive theological article for a serious general reader, able to follow careful distinctions but not presumed to know scholastic terminology. It is the first of four independent studies making up \emph{The Four Last Things}; the remaining three treat judgment, hell, and heaven, and each volume is complete in itself. Cross-references between volumes are by title only.

Its governing question is what human death is, why the tradition places it first among the last things, and where the Church's teaching about it stops. Its claim, stated compactly: the fourfold enumeration is a devotional and pedagogical ordering rather than a conciliar one, and death holds first place in it because death is the only member of the set that anyone can observe; death is defined as the separation of soul and body, a definition with dogmatic edges supplied by the teaching that the rational soul is the form of the body and is immortal, so that dying is the dissolution of one nature and not the release of a passenger; the instant of that dissolution is real and unlocatable, a result Augustine reached and the Church's sacramental law quietly presupposes; the apparent contradiction between ``death is natural'' and ``death is the wages of sin'' resolves once the two are seen to answer different questions, one about the constitution of the thing that dies and one about the withdrawal of a preserving gift; Christ's work cancels the penal meaning while leaving the natural fact standing until the resurrection, which is the exact shape of the New Testament's two tenses; what survives death is a conscious subject that is nevertheless not the person, in a condition its own principal theologian calls contrary to its nature; the Church has defined the timing and object of the vision and the timing of damnation, and almost nothing else about the interval; the accounting at death is immediate and irrevocable while the communion of the Church with the dead is not broken; the \emph{ars moriendi} is influential devotional literature of unknown authorship and no doctrinal authority whose central counsel the Catechism has retained in a single paragraph; and the Church's actual practice for the dying inverts the popular picture at both ends, since anointing is for the sick rather than only for the dying and the last sacrament is viaticum.

The argument proceeds in the order in which its own evidence becomes available: the name and the enumeration; the definition and its dogmatic supports; the unlocatable instant; the natural and penal accounts and their reconciliation; the transformation of death's meaning by Christ; the intermediate state; the judgment at death; the devotional tradition of preparation; and the Church's law and liturgy for the dying. English is the working language. No liturgical rite governs any conclusion; the canonical material is the 1983 Code for the Latin Church, as of 2026-07-26.

\subsection{Sources and their function}

\begin{threestable}{Source class}{Function in the article}{Governing boundary}
Scripture & Ecclesiasticus 7:37--40 supplies the name and its setting; Genesis 2:17 and 3:19, Wisdom 1:13 and 2:23--24, and Romans 5:12 and 6:23 the penal account; Job 14, Psalm 6:6, Ecclesiastes 12:7 and Wisdom 3:1--3 the Old Testament's development; Luke 23:43, Luke 16:22--26, 2 Corinthians 5:1--10, Philippians 1:21--24 and Apocalypse 6:9--11 the intermediate state; Hebrews 2:14--15, 9:27--28, 1 Corinthians 15:21--26 and 2 Timothy 1:10 the two tenses and the judgment; James 5:14--15 the sacrament. & Quoted from the public-domain Douay--Rheims (Challoner) in Vulgate numbering, read in the source library's tracked verse text; readings are textual and modest; no independent exegetical apparatus is claimed, and the Greek was not examined at any locus.\\
Patristic witness & Augustine, \emph{De civitate Dei} XIII, supplies the analysis of the three deaths, the evil of the separation, the unlocatable instant, the race toward death, and the argument for why the baptized still die; Cyprian, \emph{De mortalitate}, the pastoral realism about a plague that does not discriminate. & Cited as witness and record, not as demonstration; edition and translator identified; the Augustinian material verified in a tracked transcription of the Dods 1871 printing, the Cyprianic read in a web transcription of ANF~V without collation against the Latin.\\
Thomistic framework & ST I q.~75 a.~4 and q.~89 a.~1; I-II q.~85 aa.~5--6; II-II q.~164 a.~1; III q.~50 aa.~2--3; \emph{Super I Cor.} 15 lect.~2 govern the analysis of the composite, the separated soul, the natural/penal distinction, and Christ's death. & Adopted as the article's framework where Catholic theology permits schools; labeled common or school theology, never as dogma; the \emph{Super I Cor.} passages carry an explicit \latin{reportatio} caveat.\\
Definitive magisterium & Vienne, DS 902; Lateran V, DS 1440--1441; \latin{Benedictus Deus}, DS 1000--1002; Trent, session V canons 1--2 and session XIV on extreme unction, with session XIII ch.~VI on the Eucharist for the sick. & Quoted at its own register, in Latin with labeled working glosses or in a public-domain translation; the Denzinger presentation is a finding aid to the underlying acts, which are the authorities cited.\\
Authoritative non-definitive teaching & \emph{Lumen gentium} 48--50; \emph{Gaudium et spes} 14, 18; \emph{Sacrosanctum concilium} 73--74; CDF, \emph{Recentiores episcoporum Synodi} (1979); the Catechism throughout; the 1983 Code, cann.~921--923, 1004--1007. & Cited at the weight each carries and never elevated; conciliar constitutions, a curial letter, a catechetical restatement, and disciplinary law are kept distinct, and the law is dated.\\
Devotional and historical witness & \emph{The Book of the Craft of Dying} (Comper 1917) for the \emph{ars moriendi}; the \emph{Catholic Encyclopedia} for one bibliographic fact about Denys the Carthusian. & Devotional literature of no doctrinal authority, quoted in a modernized-spelling edition at two removes from the Latin; the encyclopedia used only for a dated bibliographic statement.\\
Project synthesis & The reading of Ecclesiasticus 7:40 in its context; the observation about the Catechism's placement of death; the argument that death is first because it is the only observable last thing; the two-questions reconciliation; the reading of canon 1005 against Augustine XIII.11; the mapping of the New Testament's two tenses; the reconciliation of the unlocatable instant with CCC 1022; the diagnosis of the \emph{ars moriendi}'s Platonic drift; and the closing formulations. & Original editorial synthesis argued from the checked sources; attributed to no source and creating no doctrine; the two extended syntheses carry their strongest counterargument and a defeating finding in the body.\\
\end{threestable}

\subsection{Included and excluded scope}

Included are: the scriptural origin of the phrase \emph{novissima} and the character of the fourfold enumeration; the definition of death as separation and its supports at Vienne and Lateran V; Augustine's three deaths and his aporia about the instant; the natural and penal accounts, Trent's canon, and the Thomistic reconciliation; the transformation of death's significance in Christ, with Augustine's and Cyprian's refusals to sentimentalize it; the Old Testament's reticence and the New Testament loci on the intermediate state; the separated soul in Thomistic analysis; the John XXII controversy and the definition of \latin{Benedictus Deus}; the 1979 letter and the boundary it drew; the particular judgment's substance and name and the end of the time of merit; the \emph{ars moriendi} tradition, its authorship, contents, and reception; and the Church's sacramental and liturgical practice for the dying under current Latin law.

Excluded are: the determination of death by medical criteria, organ procurement, the withholding or withdrawal of treatment, euthanasia and assisted suicide, and every other question of applied bioethics; the theology of purgatory and the whole doctrine of the final purification; hell, its punishments, and its population; the beatific vision beyond the timing and object defined in 1336; the general judgment and eschatological chronology; the resurrection of the body beyond what the definition of death requires; limbo and the fate of the unbaptized; Christology and soteriology, presupposed throughout and nowhere developed; the theology of martyrdom; funeral rubrics, cremation, burial law, and the pastoral care of the bereaved; private revelations and reported near-death experiences; the scholastic disputes over the fixity of the separated will and the twentieth-century debates over a final option at death; and any judgment about the state, culpability, or destiny of any person.

\subsection{Material qualifications}

\begin{enumerate}[leftmargin=1.45em,itemsep=0.25em]
\item Scripture is quoted only from the public-domain Douay--Rheims (Challoner) in the electronic state registered in the source library, and always in Vulgate chapter and verse numbering; a reader comparing a modern Bible will find different numbers at Psalms, Ecclesiasticus, and 1 Thessalonians in particular. Every verse quoted was checked against the registered table of divergences between that electronic state and the American 1899 edition, and none of them diverges.
\item The Greek New Testament was not examined at any locus. The essay therefore asserts nothing about the construction underlying the Vulgate's \latin{in quo omnes peccaverunt} at Romans 5:12, nor about the punctuation of Luke 23:43, and flags both where they arise.
\item The Hebrew was not examined. The remark on \latin{morte morieris} at Genesis 2:17 concerns only the Latin and English doubling.
\item Wisdom's \latin{Deus mortem non fecit} may refer chiefly to spiritual death; the essay records the ambiguity and follows Aquinas's and the Catechism's application of the verse to bodily death while flagging that the exegesis is contestable.
\item The matter/form reconciliation of the natural and penal accounts is common theology, not defined doctrine. The essay's extended synthesis defending it states the strongest objection at full strength and names the finding that would defeat it.
\item Trent's canons were read in the promulgated text in Latin and in Waterworth's English. The conciliar acts, the schemata, and the debates were not examined; the claim that Trent's canon 1 does not define the constitution of Adam's body is a claim about the promulgated canon.
\item The Denzinger presentation used for Vienne, Lateran V, \latin{Ne super his}, and \latin{Benedictus Deus} is a dated web transcription of the Latin Enchiridion that does not name the print edition it reproduces. It is used as a convenient Latin witness to acts that are themselves the authorities; no printed Denzinger was collated, and the essay's English renderings of those passages are labeled working glosses except where the official English of CCC 1023 is quoted.
\item DS 857--858, 1304--1306, and 1820, cited by CCC 1022 for the purification, were not read at their own loci. The essay reports the Catechism's citation of them and asserts nothing about their contents.
\item The report that John XXII preached the contrary opinion in 1331--1332 rests on the description of the disputed question in his own \latin{Ne super his} and on the Catechism's citation of that act; the sermons themselves were not examined.
\item The observation that the phrase ``particular judgment'' does not occur in \latin{Ne super his} or \latin{Benedictus Deus} is bounded to those texts as read on 2026-07-26. No wider search of the magisterial corpus was performed, and nothing is claimed about where the term first appears officially.
\item The \emph{Super I Cor.} passages, including \latin{anima mea non est ego}, come from the \latin{reportatio vulgata} of chapters 11--16, a reported text. They are strong evidence of the Thomistic school and weaker evidence of Thomas's own wording than their fame implies.
\item The \emph{Craft of Dying} is quoted from Comper's 1917 modernized-spelling edition of Bodleian MS 423, itself an English translation of a lost Latin original. No manuscript was consulted and no Latin wording is asserted. The attributions the tract itself makes---to Anselm for the interrogations, to Scotus, to Isidore, to ``a certain decretal''---are reported as the tract's claims and were not verified against those authors.
\item The essay's diagnosis of a Platonizing drift in the devotional literature is a critical judgment about a genre, argued from one quoted passage. It does not accuse that literature of heresy and does not extend to the \emph{ars moriendi} tradition as a whole beyond the passage cited.
\item The canonical material is the 1983 Code of Canon Law for the Latin Church, promulgated 25 January 1983 and in force since 27 November 1983, in the Holy See's English web text as of 2026-07-26. Amendments after that date were not checked. The \emph{Codex Canonum Ecclesiarum Orientalium} was not examined, and nothing here applies to an Eastern Catholic without establishing applicability. This is a study aid, not canonical or pastoral advice; anyone facing an actual decision about the sacraments for a dying person should consult the competent pastor with the complete facts.
\item Nothing in this essay states or implies a criterion for determining that a person has died, and nothing in it bears on any medical decision.
\item No claim is made about the state, salvation, or culpability of any person, living or dead.
\end{enumerate}

\subsection{Notes, translations, rights, currentness, and review}

The numbered Notes carry exact citations, edition and translation identifications, verification dates, and qualifications that would interrupt the argument; no indispensable premise lives only in a note. Scripture is quoted from the public-domain Douay--Rheims (Challoner) as registered in the source library. The \emph{Summa theologiae} is quoted from the public-domain English translation of the Fathers of the English Dominican Province with the Latin checked in the Corpus Thomisticum web corpus at the loci identified; the \emph{Super I Cor.} passages are quoted in Latin from the same corpus with labeled working glosses. Augustine's \emph{De civitate Dei} is quoted from the public-domain Dods translation of 1871, verified in the source library's tracked transcription of that printing. Cyprian is quoted from the public-domain Ante-Nicene Fathers volume~V. Trent is quoted from Waterworth's public-domain English and the 1887 Tauchnitz Latin. Conciliar and papal Latin from Vienne, Lateran V, \latin{Ne super his}, and \latin{Benedictus Deus} is quoted from a dated web transcription of the Latin Enchiridion, with English renderings labeled as working glosses except where the official English of the Catechism is used. The Catechism, the Second Vatican Council's constitutions, the 1979 letter of the Congregation for the Doctrine of the Faith, and the 1983 Code are quoted briefly from the Holy See's official English web texts with attribution. \emph{The Book of the Craft of Dying} is quoted from Comper's 1917 edition, in the public domain in the United States. Official texts and third-party translations remain outside the project's CC~BY~4.0 grant; project-created prose, images, and organization are project content.

All online witnesses were checked on 2026-07-26, and the tracked source-library artifacts on the same date; every date in this article is the UTC calendar day. The canonical material is current as of that date. This revision received internal argumentative, source-consistency, quotation, weight-classification, rights, and production review by the authoring agent. Independent scriptural, patristic, Thomistic, dogmatic, canonical, theological, literary, and ecclesiastical review is outstanding; no imprimatur, nihil obstat, or ecclesiastical approval is claimed, and internal checking is not independent review.
